\documentclass[11pt]{article}

\usepackage[english]{babel}
\usepackage[latin1]{inputenc}
\usepackage{graphicx}
\usepackage{amssymb}
\usepackage{url}

%\usepackage{mathtools}
%\mathtoolsset{showonlyrefs}

\usepackage{color}
\usepackage{natbib}
% \usepackage[style=authoryear-comp,dashed=false]{biblatex}
% \bibliography{piecewise_lyapunov}

\textwidth 15cm
\setlength{\textheight}{1.1\textheight}
\newcommand\hi{\hspace*{\parindent}}
\newcommand\vi{\vspace{\baselineskip}}
\newcommand\lac{{\mbox{{\Huge\bf L}\hspace{-0.65em}
\raisebox{-1.2ex}{\Huge\bf A}\hspace{-1.1em}
\raisebox{-0.6ex}{\Huge\bf C~}}}}


\newcommand{\Real}{\mathbb{R}}



\pagestyle{empty}

\begin{document}

\hoffset -1.5cm \voffset .5cm


\sf

\vspace*{-2.5cm} \hspace*{-0.1cm} {\small
{\mbox{\begin{minipage}{1.5cm}
\centerline{\includegraphics[width=2cm]{manchester.jpg}}
\end{minipage}}} \hspace*{0.1cm}
\begin{minipage}{11.0cm}{\sf \large \textbf{UNIVERSITY OF MANCHESTER}}\\
{\sc School of Computer Science}
\end{minipage}
} \vspace*{0mm}

\hspace*{-.7cm} {\rule[-1ex]{15cm}{0.03cm}}

\begin{flushright}
\begin{minipage}{7.0cm}\small
{\bf Please Reply to:}\\
Dr. Cordeiro\\
University of Manchester\\
School of Computer Science\\
Kilburn Building, Manchester M13 9PL\\
{\em e-mail: lucas.cordeiro@manchester.ac.uk}
\end{minipage}
\end{flushright}

\vi
\hspace*{\fill}{\small Manchester, \today.}
\vi

\begin{flushleft}
	Ahmed Louri \\
	Editor-in-Chief \\
	Department of Electrical and Computer Engineering \\
	George Washington University \\
	Washington, DC 20051 \\
%	eic_tcXemail.gwu.edu \\
	IEEE Transactions on Computers 
    \end{flushleft}
 
\indent Dear Ahmed Louri
\vi

Please find enclosed our manuscript entitled ``Multi-core synthesis and maximum satisfiability applied to optimal sizing of solar photovoltaic systems'' by Edilson Galvao, Alessandro Trindade, and Lucas Cordeiro, which we would like to submit for publication in the journal of IEEE Transactions on Computers.

The present work describes and evaluates an automated synthesis methodology to obtain the optimal sizing of photovoltaic systems. It employs model checking, specific solvers, and multi-core features to synthesize PV systems used in rural areas or where grid extension is unfeasible.

The sizing of PV systems is usually performed by hand or with commercial specialized simulation-optimization tools, which are well known in the market. However, the exploration of all design space states is impossible with those tools, and some flaws, which can be originated from the design phase, reaching the field after the PV deployment. This can cause dissatisfaction to the PV system owners and the wrong conclusion that intermittent renewable systems are not good.
   
In our study, experimental results from seven case studies showed that only the automated synthesis could find the most meaningful and detailed equipment information for PV systems in a comparative evaluation of tools. Finally, based on the fact that only since 2019 papers are dealing with automated synthesis applied to PV systems, with excellent results, but with low performance, our study is unique and with outstanding results.

The authors claim that none of the material in the mentioned paper has been published or is under consideration for publication elsewhere.

\vi
Thank you for your time.
\vi
\indent
~~~~~~~~~~~~~~~~~~~~~~~~~~~Sincerely,\\

\begin{quote}
\begin{quote}
\begin{flushright}


%\vfill

Lucas C. Cordeiro~~~~~~~
\end{flushright}
\end{quote}
\end{quote}


 
 

\end{document} 
